% Phase 2 results (Sachin)
% Numbers come from notebooks/03_dl_evaluation.ipynb: reports/dl/results_table.md, reports/dl/dl_eval_*.csv and
% the figures reports/dl/figures/dl_eval_*.png. Shared code with Phase 1: src/eval_utils.py.

\section{Deep Learning Results}
\label{sec:dl_results}

Phase~1 left two errors (Section~\ref{sec:results}): the tuned XGBoost found no fire more than 5\,km from the
day-$t$ front, and it raised false alarms around fires that were dying down. This section asks whether the final
U-Net (depth 4, focal loss, Table~\ref{tab:dl_final}), which sees the whole patch, fixes them. Every number is
computed with the same code as Phase~1 on the same 6\,727\,699 labelled test pixels; thresholds are the ones chosen
on eval. Single-model figures use seed 0; the Phase~1 predictions are refits with the committed hyperparameters
that reproduce Table~\ref{tab:results} exactly.

\subsection{Overall performance}

Over three seeds the final U-Net reaches a test PR AUC of $0.344 \pm 0.007$ against $0.255 \pm 0.002$ for XGBoost
(Table~\ref{tab:dl_seeds}). Resampling whole fire events, its lead over XGBoost is $+0.080$ (0.058 to 0.113)
for seed 0, and the interval excludes zero for every seed (Section~\ref{sec:dl_validation}). The gain is at low false-alarm rates, where an agency
would operate: recall at 1\% FPR rises from 0.388 to 0.459 and at 5\% FPR from 0.673 to 0.783, while ROC AUC moves
only from 0.905 to 0.935 (Fig.~\ref{fig:dl_pr}). The depth-3 U-Net ($0.275 \pm 0.008$) and the CNN ($0.249 \pm 0.010$) stay within about 0.02 of
XGBoost, and the depth-4 U-Net trained with weighted BCE scores 0.337, so the step that matters is depth, which lets
a pixel near the centre see the whole patch (Section~\ref{sec:dl_theory_rf}). Feeding the XGBoost map to the U-Net
as an extra channel (0.335) or averaging the two adds nothing.

\begin{figure}[t]
\centering
\includegraphics[width=0.9\columnwidth]{dl/figures/dl_eval_pr_curves.png}
\caption{Precision-recall on test, seed 0. Solid: Phase~2; dashed: Phase~1. Dots: eval-chosen thresholds.}
\label{fig:dl_pr}
\end{figure}

\subsection{Fire far from the front}

Fire 3\,km or more from the day-$t$ fire is 43\% of all next-day fire, and Phase~1 found almost none of it.
The final U-Net finds 33\% at 3--5\,km (XGBoost 4\%), and 22\%, 16\% and 21\% at 6--10, 11--20 and over 20\,km,
where XGBoost finds none (Fig.~\ref{fig:dl_dist}a). The share of missed pixels lying beyond 5\,km falls from 64\%
to 57\%. Near the fire its recall at its own threshold is lower (0.79 against 0.96 at 1\,km), but this is a
threshold effect: ranked within each distance band, the U-Net is ahead in every band, by $1.4\times$ at 1\,km and
$3$--$7\times$ from 3 to 20\,km (Fig.~\ref{fig:dl_dist}b).

\begin{figure*}[t]
\centering
\includegraphics[width=0.9\textwidth]{dl/figures/dl_eval_recall_by_distance.png}
\caption{(a) Recall of next-day fire by distance to the nearest day-$t$ fire, at each model's eval threshold;
labels give each band's share of next-day fire. (b) PR AUC within each band; grey: no skill (band fire rate).}
\label{fig:dl_dist}
\end{figure*}

Because skill tens of kilometres from any fire is surprising, we tested four other explanations. It is not
memorised places: on the 1\,186 test patches whose location is not in the training set, PR AUC beyond 5\,km is
0.179 (no skill 0.005), higher than on the 503 seen ones (0.066). It is not fire hidden under cloud: only 8.6\% of
far fire lies near a masked day-$t$ pixel, and excluding those pixels leaves recall at 0.20. It is not a
patch-level shortcut: ranking far pixels within each patch only, the U-Net reaches a median ROC AUC of 0.97 against
0.80 for XGBoost. And it does not hang on one input: blanking any channel group at inference keeps far-region PR AUC
between 0.13 and 0.15. The worst Phase~1 misses, 15--40\,km jumps and a fire growing from 8 to 260 pixels in a day,
are still missed (Fig.~\ref{fig:dl_maps}).

\subsection{False alarms on dying fires}

False positives fall from 138\,714 to 108\,623 ($-22\%$), by 25\% in patches whose fire shrinks and by 18\% in
the others (Fig.~\ref{fig:dl_dying}a), because the predicted halo is tighter and follows the shape of the fire
(Fig.~\ref{fig:dl_maps}). The U-Net does not, however, know which fire is dying. Among pixels burning on day $t$,
it still flags 81\% of those that go out, and it separates them from those that keep burning no better than
XGBoost (ROC AUC 0.66 against 0.65; Fig.~\ref{fig:dl_dying}b). Its mean score around dying and growing fires is
the same. This error is reduced, not fixed.

\begin{figure}[t]
\centering
\includegraphics[width=\columnwidth]{dl/figures/dl_eval_dying_fires.png}
\caption{(a) False-positive pixels per patch in patches whose fire shrinks (dying) or not. (b) Share of day-$t$
burning pixels flagged as fire, split by whether they still burn on day $t+1$; legend: ROC AUC between the two.}
\label{fig:dl_dying}
\end{figure}

\begin{figure*}[t]
\centering
\includegraphics[width=\textwidth]{dl/figures/dl_eval_prob_maps.png}
\caption{The 10 test patches of Fig.~\ref{fig:maps}: fire on day $t$, true fire on day $t+1$, and the scores of
XGBoost and the final U-Net; black line: eval threshold; arrow: where the wind blows.}
\label{fig:dl_maps}
\end{figure*}

\subsection{Fire events}

Pooled pixels are dominated by a few large fires (Section~\ref{sec:events}). Scored per event, the gap is much
larger: mean PR AUC 0.506 (0.480 to 0.535) for the final U-Net against 0.187 (0.161 to 0.215) for XGBoost, a
paired gain of $+0.318$ (0.288 to 0.352), with the U-Net better in 92\% of the 240 events. The gain is largest on
small fires (0.587 against 0.157), where a per-pixel model has little neighbourhood signal. At the eval threshold
the U-Net finds fire in 98.8\% of events, against 47.9\% for XGBoost and 60.8\% for persistence.

\subsection{What the U-Net looks at, and physical plausibility}

Like XGBoost, the U-Net relies mostly on where the fire is now. Blanking the day-$t$ fire mask at inference costs
0.171 PR AUC, the same as retraining without it (0.169), and the fire channel holds 47\% of the gradient saliency;
in Phase~1 the fire-neighbourhood features held most of the SHAP mass (Fig.~\ref{fig:shap}). Vegetation (0.013)
and weather (0.009) matter a little more to the trained network than retraining without them suggests (0.010,
0.002), so a network retrained without a group partly makes up for it from the others. Population is not used.

\begin{figure}[t]
\centering
\includegraphics[width=\columnwidth]{dl/figures/dl_eval_saliency_maps.png}
\caption{Gradient of the patch's predicted fire with respect to the inputs (final U-Net). Blue outline: true
day-$t+1$ fire; arrow: where the wind blows.}
\label{fig:dl_saliency}
\end{figure}

\textbf{Physical plausibility.} Two patterns agree with how fires spread. The saliency is highest on burning
pixels and on the edges of the fire ($5.2\times$ the patch average on burning pixels, $4.4\times$ at 1\,km,
$1.9\times$ at 2\,km, Fig.~\ref{fig:dl_saliency}), so the model reads the front, not the interior. The predicted
area follows the shape of the fire instead of drawing a round halo, which is what removes false alarms around
small fires. One pattern does not: neither model pushes risk
downwind. At 5\,km the score is only 1.03 (XGBoost) and 1.01 (U-Net) times higher downwind than upwind and does not
grow with wind speed, whereas next-day fire on the training split is 1.14 to 3.1 times more frequent downwind as wind
rises from below 3 to above 8\,m/s (Fig.~\ref{fig:wind}); the U-Net's saliency is as large upwind as downwind, and
rotating the wind with the patch during training changes nothing (Section~\ref{sec:dl_aug}). The wind input is a
daily value on a 4\,km grid, almost constant over a patch, which may be too coarse to locate a downwind front.
Together with the dying fires, this points to what a single day of inputs lacks (burnt fuel, suppression, next-day
weather) rather than to the model.
